\scanpage{77}
% archive sheet number: 76
\uline{vecteurs remarquables} \inserted{\uline{de directions remarquables}} (cas car. $\ne2$)

On suppose $2$ inversible : son centre $o$ identifiant $E$ et l'espace des translations $V$. Utilisant la forme quadratique, on identifie aussi $V$ et $\check V$ -- donc $\check{\mathbb P}(V)$ (directions de droites de $E$) et $\mathbb P(V)$ (directions de plans de $E$).

\margin{$S$}
$12$ vecteurs ``\inserted{doubles} sommets'' $\boxed{S_i-\underline aS_i}$

\margin{$\Delta$}
$30$ vecteurs \deleted{\unsure{arêtes}} \inserted{arêtes} $S_i-S_j$ $(\{i,j\}\in A)$ se groupent en $15$ paires de vecteurs opp.
\margin{\uline{NB} chaque arête détermine deux vect, mais deux arêtes opposées déterminent les mêmes deux}

\margin{$\Delta'$}
$30$ vecteurs coarêtes $S_i-S_j$ $(\{i,j\}\in A')$ ce sont les $\unsure{\pm}\alpha'u$ $(u\in\Delta')$
\deleted{\unsure{les mêmes en car.} $\unsure{\Delta}$ \unsure{deux}}
\[
\boxed{\Delta'=\alpha'\Delta,\quad\Delta=\alpha\Delta'}
\]
\diagram{\begin{tikzcd}[ampersand replacement=\&]
{}\arrow[r]\&{}\\
{}\arrow[u]\arrow[r]\&{}\arrow[u]
\end{tikzcd}}
\deleted{les vecteurs}

\margin{$\left[\begin{matrix}\mathcal A\\\mathcal A'\end{matrix}\right.$}
Les \inserted{\deleted{\unsure{doubles}} $30$} milieux d'arête $\deleted{\tfrac12}(S_i+S_j-2c)$ $(\{i,j\}\in A)$

\unsure{---} \inserted{$30$} de coarête $\frac12(S_i+S_j)$ $(\{i,j\}\in A')$
\[
\begin{aligned}
\Delta&=\deleted{\unsure{\xi}}2\mathcal A'&\left\{\begin{aligned}\mathcal A&=\alpha'\mathcal A'\\\mathcal A'&=\alpha\mathcal A\end{aligned}\right.\\
\Delta'&=\deleted{\unsure{\xi}}\unsure{2}\mathcal A
\end{aligned}
\]
i.e. $\Delta$ est aussi formé des $S_i+S_j$ $(\{i,j\}\in A')$
\[
\Delta'\quad\text{---}\qquad(\{i,j\}\in A)
\]

\margin{$\Phi$}
Les \inserted{$20$ triples (\unsure{ou hexuples|ou hexagones}\ldots)} barycentres des faces $\deleted{\tfrac13}(S_i\unsure{+|-}S_j\unsure{-|+}S_h\unsure{+|-}3c$ $(\{i,j,h\}\in F)$

\margin{$\Phi'$}
Les \inserted{$20$ triples} barycentres des cofaces $(\{i,j,h\}\in F')$
\margin{$\unsure{26\text{ bis}|2\text{ triples}|2\text{ trpls}}$}
\[
\Phi'=\lambda\Phi,\quad\Phi=\lambda'\Phi'\qquad(\lambda\lambda'=\mathrm{Norm}(\lambda)=\deleted{\unsure{-}}1)\qquad\lambda=2\alpha-1
\]
\deleted{\unsure{en 20}}
($\Phi$ et $\Phi'$ sont homothétiquement les sommets des Dodécaèdres inscrits dans $\Pi$ et dans le contreicosaèdre associé $\Pi'$)

\margin{$\Psi$}
Les $20$ sommets \inserted{\uline{multipliés par $10$}} des dodécaèdres circonscrits \inserted{à $\Pi$} : intersections des plans des faces adjacentes à une face

\margin{$\Psi'$}
\unsure{it.|idem} pour $\Pi'$
\[
\Psi'=\unsure{\nu|\mu}\Psi,\quad\Psi=\nu'\Psi'\qquad\deleted{\mu\mu'=\mathrm{Norm}(\mu)=1}\quad\nu=-\lambda
\]
\[
\boxed{\Psi=\mu\Phi,\quad\Psi'=\mu'\Phi'}
\qquad\begin{aligned}
\mu&=1+2\alpha=-(1+2\alpha')\\
(\mu\mu'&=-5),\quad\mu^2=5\text{ i.e. }\mu=\sqrt5
\end{aligned}
\]
\inserted{$\mu'=-\mu$}

\margin{$\Sigma$}
Les $12$ \inserted{\unsure{quintuples|quadruples}} barycentres des pentagones inscrits aux sommets
\[
2(S_1+S_2+S_3+S_4+S_0)-2c
\]
\[
\begin{aligned}
\Sigma&=\rho S=\rho'S\qquad\rho'=-\rho\quad\text{i.e. }\deleted{\rho=\rho+\rho'c}\inserted{\mathrm{Tr}\rho=0}\\
\rho&=\mu=1+2\alpha\ldots
\end{aligned}
\]
